%----------------------------------------------------------------------------------------
%	PACKAGES AND OTHER DOCUMENT CONFIGURATIONS
%----------------------------------------------------------------------------------------

\documentclass[fontsize=12pt, paper=a4, headinclude, twoside=false, parskip=half+, pagesize=auto, numbers=noenddot, open=right]{article}

\usepackage[left= 2cm,right = 2cm, bottom = 2 cm]{geometry}
\usepackage[onehalfspacing]{setspace}

\usepackage[german]{babel}
\usepackage[utf8x]{inputenc}
\usepackage{amsmath}
\usepackage{graphicx}
\usepackage{epstopdf}
\usepackage[colorinlistoftodos]{todonotes}
\graphicspath{{images/}}
\usepackage{fancyhdr}
\usepackage{caption}
\usepackage{pdfpages}
\usepackage[export]{adjustbox}
\usepackage[T1]{fontenc}            % Ermöglicht die automatische Trennung von Worten mit Umlauten
 
% Dokumentinformationen
\usepackage[
	pdftitle={Dokumentation TaskY},
	pdfsubject={Softwaredokumentation},
	pdfauthor={Michael Müller},
	pdftex=true, 
	colorlinks=true,
 	breaklinks=true,
	citecolor=black,
	linkcolor=black,	
	menucolor=black,	
	urlcolor=black
]{hyperref}

% fussnote ganz am Ende der Seite
\usepackage[bottom]{footmisc}

% nicht einrücken nach Absatz
\setlength{\parindent}{0pt}

\begin{document}

\pagestyle{empty}
\begin{titlepage}

\newcommand{\HRule}{\rule{\linewidth}{0.5mm}} % Defines a new command for the horizontal lines, change thickness here

\center % Center everything on the page
 
%----------------------------------------------------------------------------------------
%	HEADING SECTIONS
%----------------------------------------------------------------------------------------

\LARGE Hochschule für Telekommunikation Leipzig (HfTL)\\[1.2cm] % Name of your university/college
\textsc{\Large Profilierung Netzbasierte Anwendungen}\\[0.5cm] % Major heading such as course name
\textsc{\Large Profilierung Mobile Applikationen}\\[0.5cm] % Major heading such as course name
\textsc{\large Softwaredokumentation}\\[0.5cm] % Minor heading such as course title

%----------------------------------------------------------------------------------------
%	TITLE SECTION
%----------------------------------------------------------------------------------------

\HRule \\[0.4cm]
{ \LARGE \bfseries TaskY - Cache und Notifications in mobilen Webanwendungen}\\ % Title of your document
\HRule \\[1cm]
 
%----------------------------------------------------------------------------------------
%	AUTHOR SECTION
%----------------------------------------------------------------------------------------

\Large David Howon (147102)\\[1cm] 

\Large Michael Müller (147105)\\[1cm] 

%----------------------------------------------------------------------------------------
%	DATE SECTION
%----------------------------------------------------------------------------------------

{\large Wintersemester 2016/17}\\[2cm] % Date, change the \today to a set date if you want to be precise

%----------------------------------------------------------------------------------------
%	LOGO SECTION
%----------------------------------------------------------------------------------------

\includegraphics[width=0.4\textwidth]{hftl_logo.eps}\\[1cm] % Include a department/university logo - this will require the graphicx package
 
%----------------------------------------------------------------------------------------

\vfill % Fill the rest of the page with whitespace

\end{titlepage}

% Seitenränder anpassen
\newgeometry{
    left= 2.5cm, right =3.5cm, bottom = 2.5cm
}

% ============= Kopf- und Fußzeile =============
\pagestyle{fancy}
%
\lhead{Dokumentation TaskY}
\chead{}
\rhead{\slshape \leftmark}
%%
\lfoot{
    \includegraphics[scale=0.6, valign=c]{hftl_logo_sm.png}
}
\cfoot{
{\scriptsize David Howon (147102) und Michael Müller (147105)}
}
\rfoot{\textbf{\thepage}}
%%
\renewcommand{\headrulewidth}{0.4pt}
\renewcommand{\footrulewidth}{0.0pt}

% Abstand Fussnote und Fusszeile
\setlength{\footskip}{-0.5cm} 
\setlength{\skip\footins}{0.5cm} 
% Abstand zwischen Fussnoten
\setlength{\footnotesep}{0.3cm}

% \part im Inhaltsverzeichnis nicht nummerieren
\makeatletter
\let\partbackup\l@part
\renewcommand*\l@part[2]{\partbackup{#1}{}}

%Seitennummerierung neu beginnen, Zahlen [arabic], röm.Zahlen [roman,Roman], Buchstaben [alph,Alph]
\pagenumbering{Roman}
\newpage
\pagestyle{fancy}
%Inhaltsverzeichnis
\tableofcontents

\newpage
%Seitennummerierung neu beginnen, Zahlen [arabic], röm.Zahlen [roman,Roman], Buchstaben [alph,Alph]
\pagenumbering{arabic}
% pagestyle für gesamtes Dokument aktivieren
\pagestyle{fancy}

% Vorbereitende Aufgaben
\addtocontents{toc}{\protect\setcounter{tocdepth}{2}} % 0 nur > \section zeigen
\section{Anforderungen}

\subsection{allgemeine Anforderungen}

\subsection{technische Anforderungen}

\subsection{inhaltliche Anforderungen}

\newpage

% Durchführung
\addtocontents{toc}{\protect\setcounter{tocdepth}{2}} % nur bis \subsubsection (Standard 3)
\section{Durchführung}
\subsection{Betrachtung aktueller Technologien}

\subsubsection{UI-Technologien}

Betrachtung und Bewertung unterschiedlicher Möglichkeiten, die Anforderungen an das Benutzerinterface zu erfüllen.

\subsubsection{Serverlogik und Schnittstellen}

Welche Technologien sollen auf Serverseite eingesetzt werden und warum?

\end{document}